\documentclass{article}
\usepackage{amsmath, amssymb, amsthm}
\usepackage{hyperref}
\usepackage{geometry}
\geometry{margin=1in}

\title{Erd\H{o}s Problem \#321}
\date{Last updated: 2025-08-31}
\author{Source: erdosproblems.com}

\begin{document}
\maketitle

\noindent\textbf{Status:} OPEN \quad \textbf{No prize}

\noindent\textbf{Tags:} number theory, unit fractions

\medskip

\noindent\textbf{Problem Statement:}

What is the size of the largest $A\subseteq \{1,\ldots,N\}$ such that all sums $\sum_{n\in S}\frac{1}{n}$ are distinct for $S\subseteq A$?




Let $R(N)$ be the maximal such size. Results of Bleicher and Erd\H{o}s from \cite{BlEr75} and \cite{BlEr76b} imply that\[\frac{N}{\log N}\prod_{i=3}^k\log_iN\leq R(N)\leq \frac{1}{\log 2}\log_r N\left(\frac{N}{\log N} \prod_{i=3}^r \log_iN\right),\]valid for any $k\geq 4$ with $\log_kN\geq k$ and any $r\geq 1$ with $\log_{2r}N\geq 1$. (In these bounds $\log_in$ denotes the $i$-fold iterated logarithm.) See also [320] .




References

[BlEr75] Bleicher, M. N. and Erd\H{o}s, P., The number of distinct subsums of $\sum \sb{1}\spN\,1/i$ . Math. Comp. (1975), 29-42.

[BlEr76b] Bleicher, Michael N. and Erd\H{o}s, Paul, Denominators of Egyptian fractions. II . Illinois J. Math. (1976), 598-613.

\noindent\textbf{Related OEIS sequences:} A384927, A391592


\bigskip
\noindent\small{Source: \url{https://www.erdosproblems.com/321}}
\end{document}
